\documentclass[10pt,a4paper]{amsart}
\usepackage[margin=19mm]{geometry}
\usepackage{amsmath,amssymb,mathtools}
\usepackage[scr=rsfs,cal=euler]{mathalfa}
\usepackage{xfrac}
\usepackage[unicode,hidelinks,bookmarks=false]{hyperref}
\newcommand{\ie}{i.e.\ }
\newcommand{\cf}{cf.\ }
\newcommand{\resp}{respectively\ }
\newcommand{\Subsection}{\subsection}
\providecommand{\nobreakdash}{\nobreak\hbox{-}}
\setlength{\emergencystretch}{2em}
\multlinegap0pt
\usepackage{bm}
\usepackage{bbm}
\usepackage{dsfont}
\usepackage{xspace}
\usepackage{cancel}
\usepackage{mathtools}
\usepackage[shortlabels]{enumitem}
\usepackage{amsthm}
\makeatletter
\@ifundefined{theorem}{\newtheorem{theorem}{Theorem}}{}
\@ifundefined{remark}{\newtheorem{remark}[theorem]{Remark}}{}
\makeatother
\title[On the injectivity of evaluation maps induced by polynomials]{On the injectivity of evaluation maps induced by polynomials on certain algebras}
\author{Frank Kutzschebauch, Tran Nam Son, Pham Duy Vin}
\date{Source: arXiv:2508.17570v1 (CC BY 4.0)}
\begin{document}
\maketitle

\noindent\emph{Source and licence.} On the injectivity of evaluation maps induced by polynomials on certain algebras. Version arXiv:2508.17570v1 (CC BY 4.0), distributed under the Creative Commons Attribution 4.0 International licence, as recorded in the arXiv licence metadata for this exact version and shown on the abstract page https://arxiv.org/abs/2508.17570v1. Full rights notice: licenses/arxiv-2508.17570-CC-BY-4.0.txt.\par\smallskip

\noindent\textit{Editorial scope.} Counted unit: the Theorem whose source label is \texttt{variable} in the article (source0.tex lines 175--177, source label \texttt{variable}), delivered together with the article's own complete original proof of it (source0.tex lines 179--199, one \texttt{proof} environment of 21 lines). No mathematics is written, repaired or re-derived: statement and proof are the article's own text line for line. Required context retained uncounted: remark (lines 150--157). Every definition, display and label of those lines is retained unchanged. Omitted from this package and not redistributed: 1--149, 158--174, 178--178, 200--605. Nothing inside a retained excerpt is omitted or altered: every retained body is delivered exactly as the article's lines are written, so no omission receipt of any kind accompanies this package. Citations: the retained excerpts carry no citation command at all, so no bibliography entry of the article is transcribed and no reference is redistributed. Reference adaptation: none was needed and none was made; every reference inside the delivered unit resolves inside it, so no unresolved reference or citation is produced. Printed slips kept literal: no printed word, symbol, hypothesis or number of the counted statement or of its proof was altered. Layout and class: the article's own class is \texttt{amsart} with packages including amsmath, amsxtra, amscd, amsthm, amsfonts, amssymb, mathdots, bbm, graphicx, longtable, booktabs, cite, tikz-cd, hyperref; the wrapper is the lane's pinned \texttt{amsart} editorial wrapper at 10pt on A4, which restates the theorem-like environments the retained text needs and adds the article's own macro definitions that the retained text needs (none), with extra wrapper packages (bm, bbm, dsfont, xspace, cancel, mathtools, enumitem). The delivered numbering is the unit's own. Assets: no retained span contains a figure environment, a graphics-inclusion command, a PDF-page inclusion, a table or a TikZ picture; no image file is copied into this package, the package declares no asset dependency and no image file is redistributed with it. Licence: version arXiv:2508.17570v1 of this article is distributed under the Creative Commons Attribution 4.0 International licence, as recorded in the arXiv licence metadata for this exact version and shown on the abstract page https://arxiv.org/abs/2508.17570v1. A full rights notice is delivered at licenses/arxiv-2508.17570-CC-BY-4.0.txt. Characters: every retained excerpt is pure ASCII, so the delivered engine, XeLaTeX with its default Latin Modern text font, reports no missing character.
	\begin{remark} \label{remark}
		Let $\mathcal{A}$ be an associative algebra over a field $F$ and let $f\in F[x_1,x_2,\ldots,x_m]$.
		\begin{enumerate}[\rm (i)]
			\item If $\mathrm{eva}_{f,\mathcal{A}}$ is injective, then $f$ must be nonconstant. 
			\item If $f\in F[x]$ has the form: $f=ax+b$ where $a,b\in F$ and $a\neq0$, then $\mathrm{eva}_{f,\mathcal{A}}$ is injective.
			\item If $\mathcal{A}$ is unital and $\mathrm{eva}_{f,\mathcal{A}}$ is injective, then $\mathrm{eva}_{f,F}$ is also injective.
		\end{enumerate}
	\end{remark}

	\begin{theorem}\label{variable}
		Let \( F \) be a field, and let \( m \geq 2 \) be an integer. If \( F \) is either a finite field or the field $\mathbb{R}$ of all real numbers, then for any polynomial \( f \in F[x_1, x_2, \ldots, x_m] \), the evaluation map \( \mathrm{eva}_{f,F} \) cannot be injective.
	\end{theorem}

	\begin{proof}
		First, let us consider the case where \( F \) is a finite field. In this case, \( F^m \) is a finite set. If the evaluation map \( \mathrm{eva}_{f,F} \) were injective, the number of elements in \( F^m \) would have to be less than or equal to the number of elements in \( F \), which implies \( m = 1 \). This leads to a contradiction with  the assumption that $m\geq2$, so there is no polynomial whose evaluation map is injective in this case.
		
		We now turn to the case where \( F = \mathbb{R} \).	 Assume, for contradiction, that the evaluation map $\mathrm{eva}_{f,\mathbb{R}} : \mathbb{R}^m \to \mathbb{R}$ is injective. Restrict it to the unit sphere $S^{m-1} \subseteq \mathbb{R}^m$ and denote
		\[
		g := \mathrm{eva}_{f,\mathbb{R}}|_{S^{m-1}} : S^{m-1} \to \mathbb{R}.
		\]
		Since $f$ is a polynomial, it follows that $g$ is continuous, and since $\mathrm{eva}_{f,\mathbb{R}}$ is injective, so is $g$. The domain $S^{m-1}$ is compact and connected, hence the image 
		$K := g(S^{m-1})$ is also compact and connected in $\mathbb{R}$. Thus $K$ must be either a single point or a closed interval $[a,b]$.
		
		If $K$ is a point, then $g$ is constant, contradicting injectivity (see Remark~\ref{remark}). 
		If $K=[a,b]$, then $g : S^{m-1} \to [a,b]$ is a continuous bijection from a compact space to a Hausdorff space, hence a homeomorphism. 
		Pick $y \in (a,b)$ and let $x = g^{-1}(y)$. Then
		\[
		S^{m-1}\setminus \{x\} \;\cong\; [a,b]\setminus \{y\}.
		\]
		But for $m\geq 2$, the set $S^{m-1}\setminus\{x\}$ is connected 
		(indeed, homeomorphic to $\mathbb{R}^{m-1}$), 
		whereas $[a,b]\setminus\{y\}$ is disconnected. 
		This contradiction shows that no such polynomial can exist.
	\end{proof}

\end{document}
